% !TEX root = ../main.tex
% Figure: Log plot of N_PW vs bandgap accuracy

\begin{figure}[htbp]
  \centering
  \begin{tikzpicture}
    \begin{semilogyaxis}[
      width=0.9\textwidth,
      height=0.52\textwidth,
      xlabel={$N_{\mathrm{PW}}$（平面波个数）},
      ylabel={带隙相对误差（\%）},
      xmin=5, xmax=180,
      ymin=0.08, ymax=80,
      xtick={9,25,49,81,121,169},
      ytick={0.1,0.2,0.5,1,2,5,10,20,50},
      grid=major,
      grid style={dashed,gray!30},
      tick label style={font=\small},
      label style={font=\small},
      legend style={at={(0.02,0.98)},anchor=north west,font=\small,draw=none,fill=white}
    ]
      \addplot[color=blue!70!black,mark=*,line width=1.1pt]
      coordinates {
        (9,44.8618)
        (25,5.3227)
        (49,1.0298)
        (81,0.6563)
        (121,0.2435)
        (169,0.1195)
      };
      \addlegendentry{相对误差（参考解：$N_{\mathrm{PW}}=225$）}

      \addplot[color=red!70!black,dashed,line width=1.0pt] coordinates {(5,1.0) (180,1.0)};
      \addlegendentry{1\% 工程门槛}

      \addplot[color=green!55!black,dashed,line width=1.0pt] coordinates {(5,0.2) (180,0.2)};
      \addlegendentry{0.2\% 精细门槛}
    \end{semilogyaxis}
  \end{tikzpicture}
  \caption{PWEM 截断误差随 $N_{\mathrm{PW}}$ 的对数衰减趋势（本书教学脚本数据）。误差定义为 $X$ 点带隙代理 $\Delta\Omega_{23}$ 相对高截断参考解（$N_{\mathrm{PW}}=225$）的偏差。该图直观说明“先做误差-截断曲线，再给出带图结论”这一数值卫生原则。}
  \label{fig:ch12_npw_accuracy_log}
\end{figure}
